% !TeX program = xelatex
\documentclass{beamer}
\usetheme{Madrid}

% ===== 中文支持 =====
\usepackage{xeCJK}
\usepackage{ctex}
\setCJKmainfont{STKaiti}
\setCJKsansfont{STKaiti}
\setCJKmonofont{STKaiti}


\usepackage{amsmath, amssymb, amsfonts}
\usepackage{bm}
\usepackage{graphicx}
\usepackage{hyperref}

\title[LaSalle不变集原理]{LaSalle不变集原理}
%\author{非线性系统}
\date{}
\begin{document}
	
	%\begin{frame}
	%	\titlepage
	%\end{frame}
	
	\begin{frame}{正极限集 (Positive Limit Set)}
		\begin{definition}
			设 \(x(t)\) 是系统 \(\dot{x}=f(x)\) 的一个解. 若存在序列 \(\{t_n\}\) 满足
		\begin{enumerate}
			\item \(\lim\limits_{n\to\infty} t_n \to \infty\);
			\item \(\lim\limits_{n\to\infty} x(t_n) \to p\),
		\end{enumerate}
		则称 \(p\) 为 \(x(t)\) 的一个\textbf{正极限点}.
		
		所有正极限点的集合称为 \(x(t)\) 的\textbf{正极限集}，记作 \(L^+\).
		\end{definition}
	\end{frame}
	
	\begin{frame}{不变集 (Invariant Set)}
		\begin{definition}
			\begin{itemize}
			\item 设 \(x(t)\) 是 \(\dot{x}=f(x)\) 的解. 集合 \(M\) 称为\textbf{不变集}，若
			\[
			x(0)\in M \implies x(t)\in M,\quad \forall t\in\mathbb{R}.		
			\]
			\item 若仅对 \(\forall t\ge 0\) 成立，则称 \(M\) 为\textbf{正向不变集}.
		\end{itemize}
		\end{definition}
		\begin{example}
			\begin{itemize}
				\item 平衡点和极限环都是不变集。
				\item 集合 \(\Omega_c=\{x\in\mathbb{R}^n \mid V(x)\le c,\; \dot V\le 0\}\) 是不变集。
			\end{itemize}
		\end{example}
	\end{frame}
	
	\begin{frame}{局部LaSalle 不变集原理}
		\begin{definition}
			若 \(\forall \epsilon>0,\exists T>0\) 使
			$
			\operatorname{dist}(x(t),M)<\epsilon,\ \forall t\ge T,
			$
			则称当 \(t\to\infty\) 时 \(x(t)\) \textbf{趋于} \(M\),
			其中 \(\operatorname{dist}(x,M)=\inf_{p\in M}\|p-x\|\).
		\end{definition}
		\begin{theorem}
			设 \(\Omega\subset D\) 是正向紧不变集，\(V:D\to\mathbb{R}\) 连续可微且 \(\dot V(x)\le 0\). 令
			\[
			E=\{x\in\Omega \mid \dot V(x)=0\},
			\]
			\(M\) 为 \(E\) 中的\textbf{最大不变集}，则当 \(t\to\infty\) 时始于 \(\Omega\) 的每一个解都趋于 \(M\).
		\end{theorem}
	\end{frame}
	\begin{frame}
		\begin{proof}
			首先，可以证明对任何一条由 \(\Omega\) 出发的轨线 \(x(t)\) 皆有 \(\dot{V} \to 0\).
			
			考查从一点 \(x_0 \in \Omega\) 出发的轨线，对任意 \(t \geq 0\)，\(\dot{V} \leq 0\).
			
			同时 \(V(x)\) 在 \(\Omega\) 内对 \(x\) 连续（因为它对 \(x\) 可导），它在 \(\Omega\) 内有下界.
			
			由于轨线永远在正向不变集 \(\Omega\) 内，\(V(x(t))\) 对所有 \(t \geq 0\) 有下界.
			
			由\(\Omega\) 有界闭，\(V\) 有连续偏导数且\(\dot{x}\) 连续，可知 \(\dot{V}\) 一致连续.
			
			由Barbalat 引理，$
			\dot{V}(x(t)) \to 0,$
			这表明始于 \(\Omega\) 的所有轨线都趋于$E$.
		\renewcommand{\qedsymbol}{}
		\end{proof}
		\medspace
		\par 下面利用一个引理来证明所有轨线都不会收敛于 \(E\) 的其他地方，而只会收敛于 \(E\) 中的最大不变集 \(M\).
		\begin{lemma}
			若解\(x(t)\)有界且\(\forall t\ge0,\; x(t)\in D,\)则其正极限集 \(L^+\)是非空的不变闭集, 并且当 \(t\to\infty\) 时 \(x(t)\to L^+\).
		\end{lemma}
	\end{frame}
	\begin{frame}
		\begin{proof}[证明]
			令 \(x(t)\) 始于 \(\Omega\)，则 \(x(t)\) 有界且 \(x(t)\in D\)。
			由引理，\(L^+\) 非空、不变、闭，且当 \(t\to\infty\) 时 \(x(t)\to L^+\).
			即 \(\forall p\in L^+\)，存在序列 \(t_n\to\infty\) 使 \(x(t_n)\to p\).
		
			由于 \(V\) 在闭集 \(\Omega\) 上连续且单调递减，极限
			\[
			\lim_{t\to\infty} V(x(t))=:a
			\]
			存在。
			由连续性，
			\[
			V(p)=\lim_{n\to\infty} V(x(t_n))=a,\,\, \forall p\in L^+.
			\]
			故在 \(L^+\) 上恒有 \(V(x)=a\).
			又因 \(L^+\) 是不变集，沿其轨线有
			\[
			\dot V(x)=0,
			\]
			同时注意到$M$是$E$中的最大不变集，因此
			\medskip
			\par \qquad \qquad \qquad \qquad \quad \ \ \ \qquad $
			L^+ \subset M \subset E \subset \Omega.
			$
		\end{proof}
	\end{frame}
		\begin{frame}{全局LaSalle 不变集原理}
		\begin{theorem}
			若 \(V:D\to\mathbb{R}\) 连续可微，\(\dot V\le0\)，且 \(\|x\|\to\infty\) 时 \(V(x)\to\infty\)，\(E=\{\dot V=0\}\)，\(M\) 为 \(E\) 中最大不变集，则所有解全局渐近收敛于 \(M\).
		\end{theorem}
		\begin{corollary}
			设 \(V:\mathbb{R}^n\to\mathbb{R}\) 连续可微、径向无界、正定，且
			\[
			\dot V(x)\le 0,\quad \forall x\in\mathbb{R}^n.
			\]
			若由 \(\dot V(x)=0\) 确定的最大不变集 \(E\) 仅包含 \(x=0\)，则原点\textbf{全局渐近稳定}.
		\end{corollary}
	\end{frame}
\end{document}